%% ICSE 2027 NIER Track Submission
%% Double-Anonymous.  4 pp main + 1 pp references (IEEE template).
%% Compile: pdflatex; bibtex; pdflatex; pdflatex.

\documentclass[10pt,conference]{IEEEtran}

%% ---------- packages ----------
\usepackage{cite}
\usepackage{amsmath,amssymb}
\usepackage{booktabs}
\usepackage{enumitem}
\usepackage{xcolor}
\usepackage{microtype}
\usepackage[hyphens]{url}
\usepackage{balance}
\usepackage[hidelinks]{hyperref}

\setlength{\emergencystretch}{2em}
\tolerance=2000

%% ---------- metadata (anonymized) ----------
\title{Agent Observability is Not Microservice Observability}

\author{\IEEEauthorblockN{Anonymous Author(s)}
\IEEEauthorblockA{Submission anonymized for double-blind review}}

\begin{document}

\maketitle

\begin{abstract}
Large language model (LLM) agents are failing in production in ways
fifteen years of microservice observability research cannot represent.
A coding agent on \texttt{django\_\_django-10914} issues four
near-identical search queries and exhausts its iteration budget; under
the de-facto OpenTelemetry stack the failure surfaces as eight
\texttt{INTERNAL} spans with no signal that the agent was stuck. Across
a controlled 14-fault benchmark, vanilla OpenTelemetry and the current
\texttt{gen\_ai.*} semantic conventions saturate at a fault-detection
rate of 0.429 (6 of 14); the missing eight are exactly the faults that
arise from open-ended action, model-of-the-world reasoning, and
emergent multi-agent coordination. We argue this is not a
missing-attribute problem but a metamodel mismatch --- the
\emph{Cognitive-Trace Hypothesis}: agent execution is a typed
trajectory through a cognitive state space, not a request--response
chain over RPC. From this reframing we derive five concrete research
questions for the SE community, each anchored to a deliverable
artifact. One --- estimating what an autonomous agent would have done
absent a runtime intervention --- recasts SE evaluation as causal
inference on a non-stationary stochastic policy, a class of
empirical-SE question the community has not previously had to confront.
The window to set these abstractions is short: production telemetry is
already congealing around vendor-specific platforms.
\end{abstract}

\begin{IEEEkeywords}
LLM agents, observability, distributed tracing, OpenTelemetry, semantic
conventions, fault detection, software engineering research agenda
\end{IEEEkeywords}

%% ====================================================================
\section{Introduction}
%% ====================================================================

\textbf{Vignette.}
A ReAct-style coding agent~\cite{react} attacks the SWE-bench
Lite~\cite{swebench} instance \texttt{django\_\_django-10914}. Over four
consecutive iterations it issues near-identical calls to
\texttt{search\_code} with arguments
\texttt{"FilePathField"}, \texttt{"FilePathField"},
\texttt{"FilePathField allow\_files"}, \texttt{"FilePathField"} ---
never converging, finally exhausting the eight-iteration budget
without producing a patch. On the maintainer's screen this run
renders as eight \texttt{INTERNAL} OpenTelemetry spans of
indistinguishable type. There is no signal that an agent reasoning
loop occurred. The span graph satisfies the OpenTelemetry
specification~\cite{otel_spec}; the agent silently failed.

\textbf{Why the obvious vanilla-OTel detector does not work.}
A skeptic proposes a one-liner: hash the \texttt{gen\_ai.completion}
text on each \texttt{POST /v1/chat/com\-ple\-tions} span and alarm on
repetition. The proposal fails for two structural reasons. (i)~The
completion payload is a multi-thousand-token chain-of-thought blob
fusing planning, tool-argument construction, and self-critique into one
opaque string; the \emph{step within} the chain is not addressable, so
two ``different'' completions can encode the same underlying reasoning
loop and two ``identical'' completions can encode different cognitive
states. (ii)~Consecutive \texttt{chat/completions} calls in vanilla
OTel carry no causal-edge type distinguishing ``same agent looping''
from ``two independent agents querying the same backend'': the
parent-child edge is the RPC parent, not the cognitive predecessor.
The missing constructs --- a typed cognitive predecessor edge plus a
phase-addressable span --- do not exist in the standard.

This is one trace among 84 such traces in a 112-instance study using a
recently-released open-source agent-telemetry SDK~\cite{anon_sdk_2026}.
The same SDK's controlled benchmark shows that the gap is not an
artifact of any one detector: vanilla OpenTelemetry and the
\texttt{gen\_ai.*} semantic conventions both saturate at a binary
fault-detection rate of 0.429 (6 of 14 fault classes) on a controlled
fault-injection harness~\cite{anon_taxonomy_2026}. The eight
undetectable faults --- planning failure, reasoning loop, guardrail
bypass, circular delegation, memory corruption, stale retrieval, agent
misroute, and hallucination unrelated to retrieval --- have no
representation in a metamodel whose root abstraction is a typed RPC
span.

\textbf{Thesis.}
Microservice observability assumes that an execution is a DAG of typed
RPCs with bounded action spaces, no internal state worth recording, and
a single causal call tree per request. None of those assumptions hold
for LLM agents. The problem is not a missing attribute; the metamodel
was designed in 2010 for a different class of
system~\cite{dapper,tracing_practice}. We name the reframing the
\emph{Cognitive-Trace Hypothesis}: agent observability is fundamentally
a typed trajectory through a \emph{cognitive state} space, where
\emph{cognitive state} denotes the agent's evolving model-of-the-world
--- planning context, intermediate reasoning chain, memory contents,
guardrail decisions --- updated at every internal step and not
addressable through the request--response endpoint surface that
microservice tracing assumes. The hypothesis is the third generation of
a recurring SE move: distributed tracing extended single-process
logging~\cite{dapper}; the hidden-technical-debt line extended
software-systems monitoring for ML
pipelines~\cite{sculley_tech_debt,breck_ml_test_score}. Agents are the
third inflection.

This paper does three things. (1)~It identifies three structural
mismatches between microservice tracing and agent execution
(\S\ref{sec:mismatch}). (2)~It reports one emerging empirical result
that motivates the reframing (\S\ref{sec:emerging}). (3)~It articulates
a research agenda of five concrete questions, each tied to a
deliverable artifact (\S\ref{sec:agenda}), and identifies threats to
the vision (\S\ref{sec:related},~\S\ref{sec:threats}). The deepest of
those questions --- estimating what an agent would have done absent a
runtime intervention --- recasts SE evaluation as causal inference on a
non-stationary stochastic policy, a class of empirical-SE question the
community has not previously had to confront. The audience is the SE
research community: the problem is one of abstraction, conformance, and
contracts --- not of training better models --- and the window to set
these abstractions before production telemetry congeals around
vendor-specific platforms is now.

%% ====================================================================
\section{Why Microservice Observability is Insufficient}
\label{sec:mismatch}
%% ====================================================================

\subsection{The microservice telemetry model}

Dapper~\cite{dapper}, X-Trace~\cite{xtrace}, Zipkin~\cite{zipkin},
Jaeger~\cite{jaeger}, and the OpenTelemetry consensus~\cite{otel_spec}
model an execution as a DAG of typed spans
(\texttt{SERVER}, \texttt{CLIENT}, \texttt{INTERNAL},
\texttt{PRODUCER}, \texttt{CONSUMER}) attached to fixed
semantic-convention namespaces (\texttt{http.*}, \texttt{db.*},
\texttt{rpc.*}, \texttt{messaging.*}). The graph is \emph{closed}: every
node corresponds to a deployed binary's RPC boundary. Tools downstream
(latency dashboards, error-rate SLOs, trace-graph viewers) make this
closure assumption explicit. The closure is what makes
the metamodel \emph{useful}: an SRE can map every span back to a service
they can ssh into. Pivot Tracing~\cite{datadog_aps} and the broader
AIOps literature~\cite{aiops_survey,nersession} extend this picture
with dynamic queries and learned anomaly detectors, but the underlying
metamodel is unchanged.

\subsection{Three structural mismatches}

\textbf{M1. The action space is open.}
A microservice has a fixed RPC surface; an agent's next action is
sampled from a learned distribution over tools, sub-agents, internal
reasoning steps, and memory operations. The
``\texttt{POST /v1/chat/completions}'' span at the bottom of every
agent trace conceals whether the request was a planning decomposition,
a guardrail check, a chain-of-thought continuation, or a final answer.
A \texttt{SpanKind} enum sized for transport semantics cannot represent
cognition.

\textbf{M2. Causal edges are RPC parent--child, not state propagation.}
Agent failure is dominated by state corruption that flows through a
shared memory or context window, not through an HTTP header. When an
agent writes a wrong fact to its memory store at iteration~3 and a
later sub-agent retrieves it at iteration~7, no span-link in the
existing OTel model encodes the dataflow. The \texttt{traceparent}
W3C convention preserves causal containment but not causal data
flow~\cite{otel_spec}.

\textbf{M3. Aggregations are per-endpoint, not per-decision.}
Microservice SLOs are formulated over $\langle \text{endpoint},
\text{status}, \text{latency} \rangle$ tuples. An agent fleet's
service-level concerns are over $\langle \text{task}, \text{recovery
rate}, \text{cost-per-task}, \text{intervention-trigger-rate} \rangle$
tuples with no pre-defined endpoint. Practitioners currently set agent
SLOs by analogy --- ``mean latency,'' ``error rate'' --- producing
thresholds that are meaningless on a stochastic executor.

\subsection{The current state of ``LLM observability''}

Commercial GenAI tracing stacks (LangSmith~\cite{langsmith},
Langfuse~\cite{langfuse}, OpenLLMetry~\cite{openllmetry},
AgentOps~\cite{agentops}) intercept API calls and expose
prompt/completion content. They are useful, but their underlying span
model inherits the microservice metamodel: a planning step and a
summarization call are both \texttt{POST /v1/chat/completions}. As of
the OpenTelemetry GenAI SIG draft conventions (v1.30,
\texttt{gen\_ai.*} namespace), the standardized vocabulary covers
\texttt{gen\_ai.system}, \texttt{gen\_ai.request.model},
\texttt{gen\_ai.completion}, \texttt{gen\_ai.usage.*}, and the operation
names \texttt{chat}, \texttt{execute\_tool},
\texttt{retrieval}~\cite{otel_genai_sig}. None of these typed kinds or
attributes addresses the cognitive predecessor edge or the
phase-within-completion addressability that the obvious-detector
demolition in \S\ref{sec:emerging} requires.
OpenInference~\cite{openinference} goes furthest, defining typed
\texttt{LLM}, \texttt{retriever}, \texttt{chain}, \texttt{agent}, and
\texttt{guardrail} spans. It is the closest analogue to the
proposal we sketch; the gap it leaves is, however, exactly the gap
that matters: there is no typed distinction between a planning step
and an interior reasoning step (both fall under \texttt{chain}); the
\texttt{agent} span carries no source/target identifiers for typed
inter-agent \texttt{DELEGATION}; and there is no \texttt{MEMORY} span
kind for the read/write operations that drive most state-corruption
faults. The OpenInference conventions are also not part of the
OpenTelemetry standard. The OpenTelemetry GenAI SIG itself
acknowledges fragmentation and an incomplete vocabulary as open
issues~\cite{otel_genai_sig}. Recent failure taxonomies for multi-agent
systems characterize \emph{what} fails but supply no runtime telemetry
primitives capable of detecting those failures from standardized
traces: MAST~\cite{mast} catalogs 14 failure modes from manual
transcript analysis but does not define detection signals;
AgentDebug~\cite{agentdebug} learns from failure transcripts but
assumes a richer trajectory representation than vanilla OTel exposes;
Aegis~\cite{aegis} targets agent--environment mismatches but is
content-oriented, not span-typed; AgentRx~\cite{agentrx} diagnoses
post-hoc from trajectories rather than from in-flight typed spans.
AGDebugger~\cite{agdebugger} is a debugger that consumes such traces;
without typed spans it cannot attribute the failure to a cognitive
phase. Enforcement tools (NeMo Guardrails~\cite{nemo_guardrails},
GuardAgent~\cite{guardagent}) are complementary --- enforcement without
typed monitoring leaves a blind spot.

%% ====================================================================
\section{Emerging Result}
\label{sec:emerging}
%% ====================================================================

We summarize \emph{one} number that motivates the agenda; the
underlying benchmark is reported in full
elsewhere~\cite{anon_taxonomy_2026}. A 9-kind agent-telemetry SDK
\cite{anon_sdk_2026} was deployed across seven agent frameworks on 112
SWE-bench Lite instances under a ReAct coding-agent
architecture~\cite{react}. The SDK records nine span kinds covering
the orchestration phases absent from the OpenTelemetry GenAI
conventions: \texttt{AGENT}, \texttt{LLM\_CALL}, \texttt{TOOL\_CALL},
\texttt{PLANNING}, \texttt{REASONING}, \texttt{RETRIEVAL},
\texttt{GUARD\_RAIL}, \texttt{DELEGATION}, \texttt{MEMORY}.

Of 112 runs, 84 (75\%) exhausted the iteration budget without
producing a patch. All 84 exhibit a single structural pattern:
three or more consecutive
\texttt{REASONING $\to$ LLM\_CALL $\to$ TOOL\_CALL} cycles with
near-identical tool arguments. The pattern is \emph{only visible} once
the trace carries a typed \texttt{REASONING} span kind; under vanilla
OpenTelemetry the same 84 traces appear as approximately 28
undifferentiated \texttt{INTERNAL} spans each (3{,}060 spans total,
no signal). A controlled fault-injection benchmark spanning 14 fault
classes confirms the gap is structural rather than threshold-tunable:
vanilla OpenTelemetry and the current \texttt{gen\_ai.*} conventions
both saturate at a binary fault-detection rate of 0.429 (6 of 14);
the eight undetectable faults are exactly those whose detection
signal requires a span kind \emph{not present} in the existing
standard, so no detector built on the standard vocabulary in
\cite{anon_taxonomy_2026} was able to bridge the gap, independent
of threshold choice. The number we draw on for this paper is the
75\% silent-loop rate --- not the full benchmark --- because it is
the minimum evidence the vision requires: \emph{the dominant failure
mode of a deployed coding agent is silent under the de-facto
observability standard, and the silence is metamodel-level, not
detector-level.}

%% ====================================================================
\section{A Research Agenda}
\label{sec:agenda}
%% ====================================================================

Five open research questions, each anchored to a concrete artifact
the SE community can build:

\textbf{RQ1: Under what conditions is a span-kind vocabulary
\emph{saturated}, and what kinds of new agent architectures break a
given vocabulary?}
The OpenTelemetry GenAI conventions today expose nine \emph{operation
names} (\texttt{chat}, \texttt{execute\_tool}, \texttt{retrieval},
$\ldots$) but no typed \emph{kinds} for the cognitive
orchestration phases. The pilot taxonomy~\cite{anon_taxonomy_2026}
derives nine kinds from seven open-source frameworks and reports
saturation after five frameworks --- but the question of how
\emph{principled} that saturation is remains open. Does a
debate / tree-of-thought / self-consistency architecture require
additional kinds (e.g., \texttt{CRITIC}, \texttt{VOTE})? Should
memory be one kind or one per backend (vector, KV, episodic)? How do
we test, in advance, whether a candidate vocabulary will hold for
the next architecture? \emph{Artifact:} a community-curated taxonomy
with explicit saturation criteria, shipped as an OpenTelemetry
semantic-convention PR. \emph{Tractable start:} a multi-coder open-
coding study across $\geq 10$ open-source agent frameworks with
inter-coder reliability reporting, plus a stress test against
emerging multi-agent architectures.

\textbf{RQ2: How do we specify, at design time, what a correct agent
trace should look like?}
Microservice OTel has no notion of an \emph{expected} trace; the
metamodel is descriptive. For stochastic agents we need a
contract-style language --- analogous to OCL for UML or
trace-temporal-logic for distributed systems --- typed over cognitive
span kinds. Figure~\ref{fig:spec} sketches the form. The contract is
not "this trace happened" but "every trace must look like this." This
is software contracts re-imagined for stochastic executors.
\emph{Artifact:} an expected-trace specification language, a
reference checker, and a corpus of specifications mined from
successful traces. \emph{Tractable start:} mine
\texttt{must-precede}, \texttt{must-contain}, and
\texttt{must-not-repeat} patterns from the 28~successful SWE-bench
runs (the complement of the 84 silent failures in~\S\ref{sec:emerging})
in~\cite{anon_sdk_2026} and adapt LTL$_f$ for cognitive spans.

\begin{figure}[t]
\small
\hrule\vspace{2pt}
\begin{verbatim}
spec ReActAgent {
  every PLANNING span p =>
    p :precedes (TOOL_CALL | DELEGATION)
              :before LLM_CALL;
  not exists consecutive REASONING r1, r2
    where hash(r1.tool_payload)
       == hash(r2.tool_payload);
  every MEMORY.write w =>
    exists MEMORY.read r in trace
    where r.key == w.key
      and r :after w;
}
\end{verbatim}
\vspace{-4pt}\hrule
\caption{Sketch of an expected-trace specification (RQ2) over
cognitive span kinds. Clause~1 catches \emph{planning-failure}
faults; clause~2 catches the \emph{reasoning-loop} pattern from the
emerging result (\S\ref{sec:emerging}); clause~3 catches
\emph{memory-corruption} / orphan-write faults. Operators
(\texttt{:precedes}, \texttt{:before}, \texttt{:after},
\texttt{exists~in~trace}) carry standard LTL$_f$-over-event-traces
semantics typed over cognitive span kinds; full formalization is
deferred to the artifact (RQ2). The language is intentionally open ---
a research artifact, not a finished product.}
\label{fig:spec}
\end{figure}

\textbf{RQ3: How do we detect agent faults that are emergent --- visible
only in the trace as a whole, not in any single span?}
Reasoning loops, circular delegation, planning regression, and
Byzantine sub-agents are trace-level phenomena. The pilot benchmark
shows that an attribute-level detector cannot catch them even with
the right schema; one needs trace-graph queries with cognitive
typing. \emph{Artifact:} a benchmark of trace-level fault detectors
with ground-truth labels for at least the eight currently-undetectable
faults, plus precision/recall reporting against a real-LLM trace
corpus. \emph{Tractable start:} extend the MAST behavioral
taxonomy~\cite{mast} with trace-level detection signals --- the
12-of-14 MAST coverage already reported~\cite{anon_taxonomy_2026} is a
starting point, not an end state.

\textbf{RQ4: What is the right family of stochastic-SLO formalisms
for non-stationary agent fleets, and how do we express them in
telemetry?}
Latency and error-rate are the wrong primitives: they were designed
for stationary request--response systems with bounded action spaces.
Candidate primitives include task-recovery rate after intervention,
cost-per-successful-task, plan-regression-rate (the fraction of plans
superseded within $k$ steps), and the intervention-trigger-rate of
the safety policy --- but the deeper open question is whether agent
SLOs need a fundamentally different formalism (e.g., per-trajectory
percentiles, drift-aware burn-rate windows) because the underlying
agent policy is non-stationary across model updates and tool
upgrades.
\emph{Artifact:} a catalog of agent-SLO patterns with PromQL/OTel-
compatible queries, validated against at least one production agent
fleet's traces, paired with a formal characterization of which
patterns are robust under policy drift. \emph{Tractable start:}
formalize the four metrics above for the pilot corpus and report
their stability across model tiers.

\textbf{RQ5: How do we close the loop --- from observability to runtime
intervention --- without violating the principle that telemetry is
side-effect-free?}
Production agent platforms increasingly do this in an ad-hoc way:
detect a cost explosion, switch to a cheaper model; detect a
reasoning loop, inject a strategy-change prompt. The pilot
work~\cite{anon_taxonomy_2026} reports a small $+12.5$\,pp recovery
gain from one such intervention at $n{=}24$ (not statistically
significant; significance is reserved for later work). Treating
span processors as first-class \emph{actuators} re-opens questions
SE thought it had solved: how to record the intervention as a
causally-linked event~\cite{xtrace}; how to prevent intervention
loops; and --- the deepest of the three --- how to estimate the
counterfactual: what would the agent have done absent the
intervention? That is, in effect, a causal-inference problem on a
non-stationary stochastic policy, a class of empirical-software-
engineering question the community has not previously had to
confront. \emph{Artifact:} a research agenda for intervention-aware
tracing, with formal semantics for actuator-spans and a causal-
inference methodology for evaluating closed-loop interventions in
non-stationary agent populations.

Together these five questions cut across at least four core ICSE
research areas (observability, testing, formal methods, debugging)
and form a research program. None of them require advances in model
training; all of them are squarely in the SE community's
purview~--- abstractions, conformance, and contracts.

%% ====================================================================
\section{Related Vision Work}
\label{sec:related}
%% ====================================================================

Position papers on AI for SE and SE for AI have laid groundwork on
quality assurance for ML systems, on \emph{using} LLMs as SE
tools~\cite{lo_se_in_age_of_ai}, and on the methodological threats of
LLM-augmented research~\cite{nguyen_breaking_silence}; recent
visions of AI-native software engineering call for re-imagining the SE
lifecycle around foundation models~\cite{hassan_software_30}. The
structurally closest precedent is the \emph{ML-systems observability}
line: Sculley et~al.'s hidden-technical-debt
argument~\cite{sculley_tech_debt} and the ML Test
Score~\cite{breck_ml_test_score} made the case a decade ago that ML
pipelines exhibit qualitatively new observability requirements
(training--serving skew, feedback loops, entanglement). The
Cognitive-Trace Hypothesis is the third generation of that move:
distributed tracing extended single-process
logging~\cite{dapper,tracing_practice}; ML observability extended
distributed tracing; agent observability extends ML observability with
open-ended action spaces and intra-trace cognitive-state evolution.

Within the agent ecosystem, failure taxonomies
\cite{mast,agentdebug,aegis} and debugging tools
\cite{agdebugger,agentrx} make the empirical case that agent failures
are dominated by cognitive and coordination pathologies; they do not
position the gap as a metamodel problem for the SE community. The
SDK~\cite{anon_sdk_2026} on which our emerging result rests proposes
typed cognitive span kinds, but presents them as a benchmark
contribution, not as the seed of a research agenda. The contribution
of the present paper is the agenda framing.

%% ====================================================================
\section{Threats to the Vision}
\label{sec:threats}
%% ====================================================================

\textbf{T1: A narrower vocabulary wins.}
The field may converge on an LLM-only vocabulary that collapses agent
semantics into tool-call sequences. The emerging result already shows
this is structurally insufficient: tool-call sequences alone cannot
distinguish a reasoning loop from a planned multi-step search.

\textbf{T2: Frameworks consolidate.}
If agent frameworks consolidate around a single dominant API,
cross-framework abstractions may seem unnecessary. The pilot data
shows that \emph{even within a single framework} (LangChain on the
112-instance SWE-bench corpus in~\cite{anon_sdk_2026}) the structural
mismatches (M1--M3) hold.

\textbf{T3: Vendor capture.}
Production telemetry may shift toward proprietary single-vendor
platforms (Datadog LLM, Anthropic-internal, OpenAI Traces) that
absorb the standard. SE research should pre-empt this by
establishing open, typed semantics now --- the same role the
Dapper paper~\cite{dapper} played for microservices.

\textbf{T4: Chain-of-thought unfaithfulness.}
The model's verbal reasoning chain need not be a faithful record of
its internal computation~\cite{turpin_cot_unfaithful}; if a
\texttt{REASONING} span's payload is the recorded CoT, an
unfaithful CoT yields a misleading span. The agenda does not collapse
under this threat because cognitive span kinds separate cognitive
\emph{phases} from cognitive \emph{contents}: the phase-level
structural faults that the emerging result targets (loops, missing
planning, orphan memory writes) are detectable from phase typing and
causal-edge structure alone, independent of CoT content faithfulness.
Faithfulness becomes a sub-question for the contents of typed spans
(an RQ2 specification can constrain phases without trusting CoT
verbatim), not a refutation of the metamodel.

%% ====================================================================
\section{Future Plans}
\label{sec:future}
%% ====================================================================

The agenda is multi-year. Concrete near-term steps planned by the
authors and proposed for collaboration with the community:

\begin{enumerate}[leftmargin=*,nosep]
\item \textbf{Full ICSE Research-Track benchmark} extending the 14-fault
controlled study to additional architectures (tree-of-thought,
debate, hierarchical multi-agent) and a $\geq 500$-instance real-LLM
corpus.
\item \textbf{Semantic-convention proposal} of typed cognitive span
kinds to the OpenTelemetry GenAI SIG, with a public PR (citation
withheld for double-anonymous review) and a reference implementation
across at least three agent frameworks. The lived experience of the
upstream review process is itself empirical input for RQ1.
\item \textbf{Expected-trace specification language} (RQ2), released
as a stand-alone artifact with a checker and a mined-spec corpus.
\item \textbf{Controlled human study} ($n \approx 30$ SE practitioners)
measuring time-to-diagnose with versus without typed cognitive
spans, with IRB-equivalent informed consent and pre-registered
analysis.
\item \textbf{Agent-SLO pattern catalog} (RQ4), validated against a
production agent fleet's telemetry, released as an OpenTelemetry
recipe collection.
\end{enumerate}

Anonymized artifacts (the SDK source, the SWE-bench trace corpus,
the fault-injection harness, the semantic-convention specification)
will be released with the camera-ready version under an open
license.

%% ====================================================================
\section{Conclusion}
%% ====================================================================

Microservice observability is the wrong frame for autonomous agent
execution: the metamodel was designed in 2010 for typed RPC over fixed
action spaces. The Cognitive-Trace Hypothesis recasts agent
observability as typed trajectories through a cognitive state space;
the 75\% silent-loop rate under vanilla OpenTelemetry establishes the
gap as structural. Five research questions, each anchored to a
concrete artifact, define a tractable agenda --- one (RQ5) opens
causal inference on non-stationary stochastic policies as a new class
of empirical-SE question. The window to set these abstractions is
short.

%% ====================================================================
%% References
%% ====================================================================
\balance
\bibliographystyle{IEEEtran}
\bibliography{refs}

\end{document}
